% 项目科普介绍(面向非金融背景读者) —— 编译方式:XeLaTeX
\documentclass[11pt]{ctexart}
\usepackage[a4paper, margin=2.5cm]{geometry}
\usepackage{booktabs}
\usepackage{enumitem}
\setlist{nosep}
\usepackage{parskip}

\title{\textbf{这个项目是做什么的?}\\
\large ——「可选择期权」定价项目的通俗介绍}
\author{佟铸城 (Zhucheng Tong)}
\date{2026 年 7 月}

\begin{document}
\maketitle

\section{一句话概括}

这个项目要给一种特殊的金融合约——\textbf{可选择期权(chooser option)}——
算出一个公道的价格,并把算价方法做成一个人人可用的小工具。

\section{先从「定金」说起:什么是期权}

假设你看中一套房子,现价 500 万。你付 10 万定金,锁定「三个月内有权按
500 万买入」。三个月后房价涨到 550 万,你行使权利,赚了;房价跌到 450 万,
你放弃定金,损失封顶 10 万。这张「定金合同」就是金融里的\textbf{期权}:

\begin{itemize}
  \item 赌价格会\textbf{涨}的合同,叫\textbf{看涨期权(call)};
  \item 赌价格会\textbf{跌}的合同,叫\textbf{看跌期权(put)}。
\end{itemize}

买期权的妙处在于:\textbf{亏损有上限(最多亏掉「定金」),盈利没上限}。
而核心问题永远只有一个——这份「定金」收多少钱才公道?收贵了没人买,
收便宜了卖方吃亏。给这个问题算出答案,就叫\textbf{期权定价}。

\section{「可选择期权」:先上车,后选方向}

普通期权在买的那一刻就要定死方向:赌涨还是赌跌。但现实中常有这种局面:
「下个月公司发财报,股价肯定会大动,但涨是跌真说不准。」

\textbf{可选择期权}就是为此设计的:你今天付钱买入,但\textbf{可以等到约定
的「选择日」再决定}它是看涨还是看跌。相当于先付费「上车」,方向留到
消息更明朗时再选——像一张「可改签方向」的车票。

当然,这份「晚点再选」的自由不是白给的,它比普通期权更贵。\textbf{贵多少
才合理},正是本项目要精确回答的问题。

\section{价格怎么算?难点在哪?}

学术界有个经典公式(Black--Scholes 模型,1997 年诺贝尔经济学奖成果),
输入五个量就能吐出期权价格:现在的股价、约定的买卖价、离到期多久、
利率,以及\textbf{波动率}——即股价上蹿下跳的剧烈程度。

前四个都能直接查到,唯独波动率说的是\textbf{未来},只能估计。经典公式
还有个明显软肋:它假设波动率\textbf{恒定不变}。可现实里,市场平静时波动
小,金融危机、疫情、加息时波动率能翻好几倍——用「恒定」假设算出的价格,
在最需要准确的动荡时期反而最不准。

\section{我们的做法:真实数据 + 机器学习}

本项目以\textbf{摩根大通(JPM)股票}为标的,分三步改进:

\begin{enumerate}
  \item \textbf{用真实数据打底}:采集 2018--2024 共七年的股价、利率、
        市场恐慌指数(VIX)等数据——这期间恰好包含疫情崩盘、激进加息等
        极端行情,是检验模型的好考场;
  \item \textbf{先复现经典,再暴露短板}:按传统公式定价,用大量随机模拟
        交叉验证,并定量记录它在动荡时期偏差有多大;
  \item \textbf{机器学习补短板}:训练模型从历史数据中\textbf{预测}未来
        波动率(而非假设不变),再喂给定价公式;同时尝试让机器直接学习
        定价。最后用统一标准比较:新方法比经典方法准了多少。
\end{enumerate}

\section{最终交付什么?}

一个\textbf{网页版定价小工具}:输入股票代码和期权条款,它自动抓取最新
市场数据,同时给出「经典公式价」和「机器学习修正价」及其差异,并附上
一份完整的研究报告。整个项目为期八周,前两周搭数据管道,中间四周做模型,
最后两周做工具和报告。

\section*{附:八周路线图}

\begin{center}
\small
\begin{tabular}{ll}
\toprule
\textbf{阶段} & \textbf{做什么} \\
\midrule
第 1--2 周 & 采集并清洗七年市场数据,搭好自动化数据管道 \\
第 3--4 周 & 复现经典定价模型,用模拟验证,记录其失效场景 \\
第 5--6 周 & 训练机器学习模型,与经典模型正面比较 \\
第 7--8 周 & 开发网页定价工具,撰写最终研究报告 \\
\bottomrule
\end{tabular}
\end{center}

\end{document}
